\documentclass[11pt]{amsart}

\usepackage[T1]{fontenc}
\usepackage{lmodern}
\usepackage{microtype}
\usepackage{mathtools,amssymb,amsthm}
\usepackage[hidelinks]{hyperref}

\hypersetup{
  pdftitle={Pebbling numbers of four cubic 12-vertex graphs of girth 3 and
            diameter 3: one is Class 0, three are 14}
}

\newtheorem{theorem}{Theorem}
\newtheorem{lemma}[theorem]{Lemma}
\newtheorem{proposition}[theorem]{Proposition}
\theoremstyle{remark}
\newtheorem*{remark}{Remark}

\newcommand{\HoG}[1]{\#\mathrm{#1}}
\newcommand{\Reach}{\operatorname{Reach}}

\title[Four cubic $12$-vertex graphs: one Class~$0$, three with $\pi=14$]
{Pebbling numbers of four cubic $12$-vertex graphs of girth~$3$ and
 diameter~$3$: one is Class~$0$, the other three are~$14$}

\date{}

\subjclass[2020]{05C57, 05C12, 05C30}
\keywords{graph pebbling, pebbling number, Class 0 graph, cubic graph,
truncated tetrahedron, flower snark, House of Graphs, exhaustive search}

\begin{document}

\begin{abstract}
In the final remarks of \emph{The only Class~$0$ Flower snark is the smallest}
(arXiv:2505.22941v2), Bridi, Martins, Marquezino and de Figueiredo suggest
investigating whether the pebbling number of the House of Graphs graphs
$\HoG{1395}$ (the truncated tetrahedron), $\HoG{44170}$ and $\HoG{44172}$ is
indeed $13$, and whether $\HoG{6698}$ is Class~$0$. We settle both questions for
the four explicit cubic $12$-vertex graphs of girth~$3$ and diameter~$3$ whose
adjacency lists are printed below; these lists were served by the House of Graphs
REST API for those four ids, but we do not verify that identification, and every
statement here is a statement about the lists as printed. For three of them the
suggested value $13$ does not hold: $\pi=14$, and each excess is witnessed by an
explicit $13$-pebble unsolvable configuration whose unsolvability is certified by
the finite move-closed set of configurations reachable from it, so that lower
bound rests on the closure of three single configurations rather than on an
exhaustive sweep of the configurations of a given total. The fourth has
$\pi=12$, i.e.\ it \emph{is} Class~$0$, and it is not isomorphic to $J_3$; so the
Flower snark $J_3$ is not the only Class~$0$ graph among the cubic $12$-vertex
graphs of girth~$3$ and diameter~$3$.
\end{abstract}

\maketitle

\section{Introduction}

We use the definitions of the source paper verbatim \cite[\S2]{BMMF}. A
\emph{configuration} on a graph $G$ is a map $C\colon V(G)\to\mathbb{N}$, where
$C(v)$ is the number of pebbles at $v$; its \emph{total} is $\sum_v C(v)$. A
\emph{pebbling move} removes two pebbles from a vertex $v$ and places one pebble
on a neighbour $u\in N(v)$. If some sequence of pebbling moves places a pebble on
a target vertex $r$, then $C$ is \emph{$r$-solvable}; otherwise $C$ is
\emph{$r$-unsolvable}. The \emph{pebbling number} $\pi(G)$ is the least $t$ such
that every configuration of total $t$ is $r$-solvable for every $r\in V(G)$.
Always $\pi(G)\ge n(G)$, and $G$ is \emph{Class~$0$} if $\pi(G)=n(G)$.

The source paper proves that the Flower snark $J_3$ is Class~$0$, i.e.\
$\pi(J_3)=12$, and closes with the following suggestion, which we quote verbatim
from the final remarks of its \texttt{arXiv} v2 \LaTeX{} source:

\begin{quote}\small
As a suggestion of application, one natural direction is to consider other cubic
graphs with $12$ vertices, girth~$3$, and diameter~$3$---the same parameters as
the Flower snark $J_3$. According to the House of Graphs
database~\cite{HoG1}, only five such graphs exist. Besides $J_3$, the others
include graphs with IDs \#1395, \#6698, \#44170, and \#44172. The graph \#1395 is
known as the truncated tetrahedral graph. We find unsolvable configurations with
$12$ pebbles for the graphs \#1395, \#44170, and \#44172 (see
Figure~[reference omitted]), implying that these graphs are not Class~$0$. One
could investigate whether the pebbling number of these three graphs is indeed
$13$; and whether graph \#6698 is Class~$0$, i.e., determine whether $J_3$ is the
only Class~$0$ graph among the five with the same parameters.
\end{quote}

Only the presentation is altered: the citation of the database is re-keyed to our
bibliography (our \cite{HoG1} is the source's own database reference), the
bracketed note stands for the source's cross-reference to its figure of
$12$-pebble unsolvable configurations, which we do not reproduce, and a footnote
giving the database address is omitted.

The suggestion has two halves. The first offers a value tentatively rather than
asserting it: the source's own figure of $12$-pebble unsolvable configurations
already gives $\pi\ge13$ for the three graphs, so what its suggested value~$13$
would require is the upper bound $\pi\le13$. The second half is posed neutrally.

Throughout, $\HoG{1395}$, $\HoG{6698}$, $\HoG{44170}$ and $\HoG{44172}$ are names
for the four adjacency lists printed in \S\ref{sec:graphs}, and nothing below is a
claim about the House of Graphs database. Those lists were served by the
database's REST API for the four ids named in the quotation; we could not
afterwards read the database's holdings, since its search endpoint returned
HTTP~401 to us, so the identification of the printed lists with those ids is not
verified here and is not verified by the accompanying program. Three of the lists
are, however, pinned to the source's own figure of $12$-pebble unsolvable
configurations by an exact edge-set identity: reading the node labels of each of
its three subfigures as $\text{API index}=\text{label}-1$ reproduces the fetched
adjacency list of the corresponding graph exactly, $18=18$ edges with empty
symmetric difference. For $\HoG{6698}$, which the source's figure does not show,
there is no such pin.

\begin{theorem}\label{thm:13}
$\pi(\HoG{1395})\ge14$, $\pi(\HoG{44170})\ge14$ and $\pi(\HoG{44172})\ge14$. In
particular the suggested value $13$ does not hold for any of the three graphs.
\end{theorem}

Theorem~\ref{thm:13} is carried by the three configurations of
Table~\ref{tab:witnesses}, each of total $13$; \S\ref{sec:certs} certifies each
of them by the finite set of configurations reachable from it, which is closed
under pebbling moves and contains no configuration with a pebble on the target.
Checking such a certificate is the deterministic closure of one configuration:
still a finite search of the configurations reachable from it, but not an
enumeration of the configurations of a given total. The three closures
have $6358$, $1912$ and $1711$ members; we do not print them, so at that size the
check is delegated to a program.

\begin{theorem}\label{thm:exact}
$\pi(\HoG{1395})=\pi(\HoG{44170})=\pi(\HoG{44172})=14$ and
$\pi(\HoG{6698})=12$. In particular $\HoG{6698}$ is Class~$0$, so $J_3$ is
\emph{not} the only Class~$0$ cubic $12$-vertex graph of girth~$3$ and
diameter~$3$.
\end{theorem}

The upper bounds in Theorem~\ref{thm:exact} rest on exhaustive search, organised
by the two elementary lemmas of \S\ref{sec:exhaustive}; the lower bounds are
certificates. Nothing in the source paper is contradicted: both of its proved
statements, $\pi(J_3)=12$ and $\pi\ge13$ for the three graphs, are confirmed
here.

\section{The graphs}\label{sec:graphs}

Vertices are $0,\dots,11$; the lists below are the adjacency lists served by the
House of Graphs REST API for the four ids, under
$\text{API index}=\text{HoG label}-1$. Every statement in this note about the five
named graphs is a statement about these lists as printed, so a reader who
distrusts the database can still check it.

\medskip
\noindent\texttt{\#1395} (the truncated tetrahedral graph):
\[\begin{array}{llll}
 0\colon 1,2,3 & 1\colon 0,8,9 & 2\colon 0,3,11 & 3\colon 0,2,10\\
 4\colon 5,6,10 & 5\colon 4,7,11 & 6\colon 4,8,10 & 7\colon 5,9,11\\
 8\colon 1,6,9 & 9\colon 1,7,8 & 10\colon 3,4,6 & 11\colon 2,5,7
\end{array}\]

\noindent\texttt{\#6698} (the graph the source asks about):
\[\begin{array}{llll}
 0\colon 1,2,3 & 1\colon 0,8,9 & 2\colon 0,3,11 & 3\colon 0,2,10\\
 4\colon 5,8,10 & 5\colon 4,8,11 & 6\colon 7,9,11 & 7\colon 6,9,10\\
 8\colon 1,4,5 & 9\colon 1,6,7 & 10\colon 3,4,7 & 11\colon 2,5,6
\end{array}\]

\noindent\texttt{\#44170}:
\[\begin{array}{llll}
 0\colon 1,2,3 & 1\colon 0,10,11 & 2\colon 0,10,11 & 3\colon 0,8,9\\
 4\colon 5,6,8 & 5\colon 4,7,8 & 6\colon 4,9,10 & 7\colon 5,9,11\\
 8\colon 3,4,5 & 9\colon 3,6,7 & 10\colon 1,2,6 & 11\colon 1,2,7
\end{array}\]

\noindent\texttt{\#44172}:
\[\begin{array}{llll}
 0\colon 1,2,3 & 1\colon 0,9,11 & 2\colon 0,10,11 & 3\colon 0,7,8\\
 4\colon 5,6,8 & 5\colon 4,9,10 & 6\colon 4,9,10 & 7\colon 3,8,11\\
 8\colon 3,4,7 & 9\colon 1,5,6 & 10\colon 2,5,6 & 11\colon 1,2,7
\end{array}\]

\noindent The control $J_3$ is taken from the source's own figure for $J_3$ with
its labels $A,\dots,L$ read as $0,\dots,11$, i.e.\ the edges
$AB$, $AC$, $AD$, $BE$, $BF$, $CG$, $CH$, $DI$, $DJ$, $EK$, $FL$, $EF$, $GK$,
$IK$, $HL$, $JL$, $HI$, $GJ$.

\begin{proposition}\label{prop:objects}
Each of the five graphs is a simple connected cubic graph on $12$ vertices with
$18$ edges, girth~$3$, radius~$3$ and diameter~$3$; they are pairwise
non-isomorphic. Moreover $\HoG{1395}$ is isomorphic to the truncation of $K_4$,
and the graph just read off the source's figure is isomorphic to the standard
Flower snark $J_3$ (three stars $K_{1,3}$ with centres $a_i$ and leaves
$b_i,c_i,d_i$; a triangle on the $b_i$; a $6$-cycle $c_0c_1c_2d_0d_1d_2$).
\end{proposition}

Every clause of Proposition~\ref{prop:objects} is recomputed from the printed
lists by the accompanying program.

\section{Three certificates}\label{sec:certs}

For a configuration $C$ let $\Reach(C)$ be the set of configurations obtainable
from $C$ by (possibly empty) sequences of pebbling moves.

\begin{lemma}\label{lem:cert}
Let $r\in V(G)$ and let $S$ be a set of configurations with $C\in S$, $S$ closed
under pebbling moves, and $D(r)=0$ for every $D\in S$. Then $C$ is
$r$-unsolvable, and $\pi(G)\ge\lvert C\rvert+1$ where $\lvert C\rvert$ is the
total of $C$.
\end{lemma}

\begin{proof}
By induction along a sequence of moves, every configuration reachable from $C$
lies in $S$, hence has no pebble on $r$; so no sequence of moves places a pebble
on $r$. As $C$ has total $\lvert C\rvert$ and is not $r$-solvable, $\pi(G)>\lvert
C\rvert$.
\end{proof}

The set $S=\Reach(C)$ is closed by construction, and it is finite because a
pebbling move strictly decreases the total. So a certificate is just the finite
list $\Reach(C)$, and checking it is bookkeeping: generate the moves, confirm
none lands a pebble on $r$.

\begin{table}[ht]
\caption{The three witnesses. Each has total $13$ and $C(r)=0$, and each
$\Reach(C)$ is closed under moves and free of pebbles on $r$.}
\label{tab:witnesses}
\[
\begin{array}{clcl r}
 & \text{graph} & r & C=(C(0),\dots,C(11)) & \lvert\Reach(C)\rvert\\\hline
W_1 & \HoG{1395}  & 0 & (0,0,0,0,3,4,1,3,0,1,1,0) & 6358\\
W_2 & \HoG{44170} & 1 & (0,0,1,1,5,5,0,0,1,0,0,0) & 1912\\
W_3 & \HoG{44172} & 5 & (3,1,1,0,0,0,1,7,0,0,0,0) & 1711
\end{array}
\]
\end{table}

\begin{proof}[Proof of Theorem~\ref{thm:13}]
Apply Lemma~\ref{lem:cert} to $S=\Reach(W_i)$ for $i=1,2,3$ with the targets of
Table~\ref{tab:witnesses}. Each $W_i$ has total $13$, so
$\pi\ge14$ for the corresponding graph.
\end{proof}

\begin{sloppypar}
The source's own three $12$-pebble configurations, read off its figure of
unsolvable configurations under $\text{API index}=\text{label}-1$, are
re-certified unsolvable in the same way, which confirms its bound $\pi\ge13$:
$\HoG{1395}$, $r=0$, $(0,0,0,0,0,7,3,0,0,1,1,0)$, $\lvert\Reach\rvert=1262$;
$\HoG{44170}$, $r=5$, $(5,1,1,0,0,0,0,0,0,0,5,0)$, $1251$; and $\HoG{44172}$,
$r=0$, $(0,0,0,0,7,1,1,0,0,1,1,1)$, $483$.
\end{sloppypar}

\section{The exact values}\label{sec:exhaustive}

Fix $G$ and a target $r$. A configuration with $C(r)\ge1$ is $r$-solvable
trivially, so let
\[
 U_s(r)=\{\,C : \textstyle\sum_v C(v)=s,\ C(r)=0,\ C\ \text{is $r$-unsolvable}\,\}.
\]
Then $\pi(G)=\max_{r}\min\{s : U_s(r)=\emptyset\}$ follows from the definition
together with Lemma~\ref{lem:down} below. Two lemmas make $U_s(r)$ computable
without ever enumerating the whole space.

\begin{lemma}[downward closure]\label{lem:down}
If $C'\le C$ pointwise and $C$ is $r$-unsolvable, then $C'$ is $r$-unsolvable.
Consequently $U_s(r)=\emptyset$ implies $U_{s'}(r)=\emptyset$ for all $s'\ge s$,
and every $C\in U_s(r)$ arises as $C'+e_v$ for some $C'\in U_{s-1}(r)$ and some
$v\ne r$.
\end{lemma}

\begin{proof}
Every move legal at $C'$ is legal at $C$, and applying the same move to both
preserves $C'\le C$; so a sequence of moves solving $C'$ solves $C$, and the
first claim is the contrapositive. If $C\in U_{s+1}(r)$, remove one pebble from
any $v\ne r$ with $C(v)\ge1$: the result lies in $U_s(r)$ by the first claim,
which gives both remaining assertions.
\end{proof}

\begin{lemma}[one-move recurrence]\label{lem:rec}
Let $C$ have total $s\ge1$ and $C(r)=0$. Then $C\in U_s(r)$ if and only if for
every $v$ with $C(v)\ge2$: $r\notin N(v)$, and $C-2e_v+e_u\in U_{s-1}(r)$ for
every $u\in N(v)$.
\end{lemma}

\begin{proof}
$C$ is $r$-solvable iff some first move leads to an $r$-solvable configuration or
places a pebble on $r$ directly; a move from $v$ to $r$ is available exactly when
$C(v)\ge2$ and $r\in N(v)$. A move lowers the total by $1$, so the recurrence is
acyclic and needs no fixed point. Configurations all of whose entries are at most
$1$ admit no move and are therefore in $U_s(r)$, which is the base case.
\end{proof}

Lemmas~\ref{lem:down} and~\ref{lem:rec} give a complete algorithm: compute
$U_0(r)=\{0\}$, and obtain $U_s(r)$ by adding one pebble in every possible way to
every member of $U_{s-1}(r)$ --- complete by Lemma~\ref{lem:down} --- and testing
each candidate by Lemma~\ref{lem:rec}. The search stops at the first empty level,
which by Lemma~\ref{lem:down} is $\min\{s:U_s(r)=\emptyset\}$.

\begin{proof}[Proof of Theorem~\ref{thm:exact}]
Running the algorithm on the printed adjacency lists for all $12$ targets of each
graph gives the level sizes $\lvert U_s(r)\rvert$ recorded in
Table~\ref{tab:counts}: for $\HoG{6698}$ (and for the control $J_3$) every target
already has $U_{12}(r)=\emptyset$, while $U_{11}(r)\ne\emptyset$, so
$\pi(\HoG{6698})=12=n$ and $\HoG{6698}$ is Class~$0$; for the other three graphs
$U_{13}(r)\ne\emptyset$ for at least one target and $U_{14}(r)=\emptyset$ for all
targets, so $\pi=14$.
\end{proof}

\begin{table}[ht]
\caption{$\lvert U_s(r)\rvert$ by target $r=0,\dots,11$, and the sums over all
$12$ targets. Every entry is a count of $r$-unsolvable configurations; a
configuration with a pebble on $r$ is $r$-solvable by definition and is not
counted.}
\label{tab:counts}
\[
\begin{array}{llr}
s & \text{per-target vector} & \text{sum}\\\hline
12 & \HoG{1395}\colon (165,165,165,165,165,165,165,165,165,165,165,165) & 1980\\
12 & \HoG{44170}\colon (0,229,229,0,91,91,0,0,91,3,0,0) & 734\\
12 & \HoG{44172}\colon (8,0,0,94,4,153,153,94,18,5,5,8) & 542\\
12 & \HoG{6698}\colon (0,0,0,0,0,0,0,0,0,0,0,0) & 0\\
12 & J_3\colon (0,0,0,0,0,0,0,0,0,0,0,0) & 0\\[2pt]
13 & \HoG{1395}\colon (2,2,2,2,2,2,2,2,2,2,2,2) & 24\\
13 & \HoG{44170}\colon (0,20,20,0,4,4,0,0,4,0,0,0) & 52\\
13 & \HoG{44172}\colon (0,0,0,0,0,6,6,0,0,0,0,0) & 12\\
13 & \HoG{6698}\colon (0,0,0,0,0,0,0,0,0,0,0,0) & 0\\[2pt]
14 & \text{all five graphs}\colon (0,\dots,0) & 0
\end{array}
\]
\end{table}

The sweep accounts for the whole space: for a fixed target the configurations with
$C(r)=0$ and total $s$ number $\binom{s+10}{10}$, which is $646\,646$ at $s=12$
and $1\,144\,066$ at $s=13$, so the Class-$0$ statement for $\HoG{6698}$ is a
statement about all $12\cdot646\,646=7\,759\,752$ of them. The three witnesses of
Table~\ref{tab:witnesses} are members of the corresponding $U_{13}$, which ties
the certificates of \S\ref{sec:certs} to the sweep.

So the answer to the source's restatement of its own question --- ``determine
whether $J_3$ is the only Class~$0$ graph among the five with the same
parameters'' --- is no: the graph printed here as $\HoG{6698}$ is Class~$0$ as
well, and it is not isomorphic to $J_3$.

\section{Scope}\label{sec:scope}

The upper bounds have no hand-checkable certificate here: that
$U_{12}=\emptyset$ for $\HoG{6698}$, and $U_{14}=\emptyset$ for the other three,
is proved by exhaustion over a search space whose completeness is itself proved
(Lemma~\ref{lem:down}), but it is exhaustion, and no weight-function or LP-dual
certificate for $\pi(\HoG{6698})\le12$, checkable with integer arithmetic and no
program, is offered. Theorem~\ref{thm:13} asks for less trust, resting on the
closure of three single configurations of a few thousand members each rather than
on an exhaustive sweep, but those closures are not printed here either, so it too
is a machine check, only a small and local one.

Only the five graphs printed in \S\ref{sec:graphs} are treated. How many cubic
$12$-vertex graphs of girth~$3$ and diameter~$3$ there are, and what their
pebbling numbers are, is not settled here; in particular the source's report that
the database holds five such graphs is neither used nor contradicted.

We give exact values, not a characterisation: nothing here explains why
$\HoG{6698}$ is Class~$0$ while $\HoG{1395}$, $\HoG{44170}$ and $\HoG{44172}$ are
not. Neighbouring cells are untouched: cubic $12$-vertex graphs of girth~$4$ or
girth~$5$ at diameter~$3$, girth-$3$ members of diameter at least~$4$, every
$n\ne12$, and the source's own question about extending its method beyond
diameter~$3$. The best published diameter-$3$ upper bound,
$\pi\le\frac{3}{2}n+O(1)$~\cite{bukh,psy}, gives roughly $20$ at $n=12$ and
cannot separate $12$ from $13$ from $14$, so no prior theorem forces any of the
values above.

\begin{remark}[verification]
The program \texttt{verify.py} accompanying this note recomputes, from the
adjacency lists of \S\ref{sec:graphs} transcribed into it, using only the Python
standard library and exact integer arithmetic, the graph parameters and
identifications of Proposition~\ref{prop:objects}, the reachability certificates
of \S\ref{sec:certs} with their sizes, all of Table~\ref{tab:counts}, and the
pebbling number of each of the five named graphs; its recorded run also reports
quantities this note does not print, among them the pebbling numbers of eleven
further cubic $12$-vertex graphs of girth~$3$ and diameter~$3$, which nothing here
claims. It does not parse this note, so the transcription of the printed objects
into its inputs is not checked by it; it does not fetch or parse any House of
Graphs record, so the identification of the four adjacency lists with the ids
$\HoG{1395}$, $\HoG{6698}$, $\HoG{44170}$ and $\HoG{44172}$ is not checked by it
either, as its own scope note records.
\end{remark}

\begin{thebibliography}{9}

\bibitem{BMMF}
G.~A. Bridi, A.~L.~A. Martins, F.~L. Marquezino, C.~M.~H. de Figueiredo,
\emph{The only Class $0$ Flower snark is the smallest},
arXiv:2505.22941v2 [math.CO], 2 June 2025.
\newblock DOI \texttt{10.48550/arXiv.2505.22941}. The quotation in \S1 is from the
final remarks of the v2 \LaTeX{} source; the definitions quoted in \S1 and the
theorem $\pi(J_3)=12$ are from its Section~2.

\bibitem{HoG1}
G.~Brinkmann, K.~Coolsaet, J.~Goedgebeur, H.~M\'elot,
\emph{House of Graphs: a database of interesting graphs},
Discrete Appl. Math. \textbf{161} (2013), 311--314.
\newblock This is the reference the source paper cites for the database.

\bibitem{HoG2}
K.~Coolsaet, S.~D'hondt, J.~Goedgebeur,
\emph{House of Graphs 2.0: A database of interesting graphs and more},
Discrete Appl. Math. \textbf{325} (2023), 97--107.
\newblock The four adjacency lists of \S\ref{sec:graphs} were served by this
version's REST API, \texttt{houseofgraphs.org/api/graphs/<id>}.

\bibitem{chung}
F.~R.~K. Chung,
\emph{Pebbling in hypercubes},
SIAM J. Discrete Math. \textbf{2} (1989), 467--472.

\bibitem{psv}
L.~Pachter, H.~S. Snevily, B.~Voxman,
\emph{On pebbling graphs},
Congr. Numer. \textbf{107} (1995), 65--80.

\bibitem{bukh}
B.~Bukh,
\emph{Maximum pebbling number of graphs of diameter three},
J. Graph Theory \textbf{52} (2006), 353--357.

\bibitem{psy}
L.~Postle, N.~Streib, C.~Yerger,
\emph{Pebbling graphs of diameter three and four},
J. Graph Theory \textbf{72} (2013), 398--417.

\bibitem{ADFHS}
M.~Adauto, C.~M.~H. de Figueiredo, G.~Hurlbert, D.~Sasaki,
\emph{On the pebbling numbers of Flower, Blanu\v{s}a and Watkins snarks},
Discrete Appl. Math. \textbf{361} (2025), 336--346.
\newblock Proves $\pi(J_3)\le13$; it does not consider any of the four graphs
above.

\end{thebibliography}

\end{document}
